% 第 3 章：详细 IEC 60909 计算
\section{详细 IEC 60909 计算}\label{sec:sc-detailed}

\subsection{序矩阵构建与选择性求逆}
详细路径 \srcpath{run\_short\_circuit\_detailed}（\srcpath{include/hacdcpf/analysis/short\_circuit.hpp:run\_short\_circuit\_detailed}）在
\srcpath{build\_sc\_admittance\_matrices} 中建稀疏序矩阵，\textbf{绝不}构造完整 $Z_{bus}$：对故障
母线 $k$ 将每个所需序矩阵分解一次并解 $Y_s z_k=e_k$，$z_k$ 为所选逆列；非故障电流指标另需
选择性逆对角元，以有界稠密 RHS 块提取并仅保留各块对角（\fld{inverse\_rhs\_batch\_size}，默认 32）。

\subsection{元件序阻抗}
\begin{itemize}
  \item \textbf{无源支路}：$z_1=r+jx$、$y_1=1/z_1$；零序在有 \fld{r0\_pu}/\fld{x0\_pu} 时用 $z_0=r_0+jx_0$，
        缺失则不提供零序路径（除非投影填回退值）；非标准变比 $t$ 以标准 pi/tap 块插入；
  \item \textbf{外部电网}：由 \fld{ikq\_ka}+\fld{x\_r}、直接标幺 $r/x$、或短路功率给出——
        $Z_{ext,\Omega}=cU_{nQ}/(\sqrt3 I_{kQ}'')$，或 $|Z_{ext,pu}|=cS_b/S_{sc}$；
  \item \textbf{同步发电机}：$z_G''=(r_a+jx_d'')S_b/S_{rG}$，先按
        $(U_{rG}/U_b)^2$ 换到母线电压基准，再乘含 $U_n/U_{rG}$ 与 $p_G$ 的 $K_G$；
        电站机组外部故障时发电机与升压变同时乘 $K_S$，机端内部故障则变压器恢复原阻抗、
        发电机改用内部故障 $K_G$；
  \item \textbf{线路}：最小短路按 IEC 60909-0:2016 式 (14) 使用
        $R_\theta=[1+\alpha(\theta_e-\theta_{ref})]R_{ref}$；最大短路保持参考温度电阻；
  \item \textbf{变压器}：两绕组支持 YNyn/YNd/Dyn/Yyn/YNy 及稳定三角绕组写法，
        不接地 Yy/Yd/Dy/Dd 阻断零序；中性点电抗、励磁零序阻抗和分配系数进入 T 型等值。
        三绕组成对短路试验阻抗先乘各自 $K_T$，变换为三支星形绕组后以含理想变比的四节点
        印章组装，再对内部星点作 Schur 消元；零序按三个绕组的接地/三角/不接地语义建立同构印章。
\end{itemize}

\subsection{特征电流}
详细结果 \fld{SCDetailedBusResult}（:SCDetailedBusResult）给出 IEC 60909 全套特征电流：
\begin{itemize}
  \item \textbf{初始对称短路电流} $I_k''$（\fld{ikss\_ka}），并分列发电机/电机/静止发电机/外网/
        换流器贡献（\fld{ikss\_gen\_contrib\_ka} 等）与去电机值 \fld{ikss\_no\_motor\_ka}；
  \item \textbf{峰值电流} $i_p$：方法 C 建立 $f_c/f=0.4$ 的独立正序稀疏网络，所有
        无源/机器电抗按频率缩放，同步发电机采用 $R_{Gf}$，由 $Z_{c,kk}$ 求单一
        $\kappa_C$，并按 IEC 60909-0 (59)
        \begin{equation}
        i_p=\sqrt2\bigl(\kappa_C I_{k1}''+I_{k2}''\bigr)
        \end{equation}
        合成，其中 $I_{k2}''$ 是静止发电机/跟网换流器限流电流源贡献。方法 A/B 同样对
        等效电压源总电流使用一个方法系数，电流源项取 $\kappa=1$；非故障母线采用单
        $\kappa$ 转移近似。基础冲击系数
        $\kappa=1.02+0.98\,e^{-3R/X}$（\srcpath{src/short\_circuit/short\_circuit.cpp:calculate\_kappa\_basic\_sc}；方法选择
        \fld{SCKappaMethod} A/B/C）。方法 A 由故障电压解筛出故障电压等级的实际参与支路和
        馈入变压器，并对电源内阻与全部电压源贡献统一使用最小 $R/X$ 所得 $\kappa_A$；方法 B
        且 \fld{Meshed} 时乘 1.15，低压上限 1.8、中高压上限 2.0；\fld{Auto} 通过独立
        电源路径数判定是否网状，并按标准对任一路径 $R/X<0.3$ 取消 1.15 修正；
  \item \textbf{开断电流} $I_b$（\fld{ib\_ka}）：$\mu$ 在电流比不大于 2 时取 1，
        在 20/50/100/250 ms 标准曲线间按 \fld{breaking\_time\_s} 插值；多馈三相故障按
        公式 (77) 使用 $|Z_iI_{ki}''|/(cU_n/\sqrt3)$ 权重，不平衡故障按公式 (78)--(80)
        取 $I_b=I_k''$；
  \item \textbf{稳态电流} $I_k$（\fld{ik\_ka}）：先以单机贡献
        $I_{kG}''>2I_{rG}$ 判定近端。单馈近端使用显式制造商/标准曲线字段
        \fld{sc\_lambda\_max/sc\_lambda\_min} 计算 $I_{kG}=\lambda I_{rG}$；缺失所需
        $\lambda$ 时返回 \fld{invalid\_generator\_data}，不由 $x_d/x_q$ 猜值。端子馈电静态励磁的
        机端故障在最大/最小计算中均取 $\lambda_{\min}$；多馈近端取去除全部异步电机后的
        $I_k=I_{bmo}$；远端网络馈入保持初始值，异步电机稳态为零；
  \item \textbf{热等效电流}：\srcpath{calculate\_thermal\_factors\_sc} 按 IEC 60909-0:2016
        附录 A 计算 $m(\kappa,fT_k)$ 与 $n(I_k''/I_k,T_k)$，含 $\kappa\to1/2$、$I_k\to0$
        极限及图 19 的 $1< I_k''/I_k <1.25$ 插值，最后
        $I_{th}=I_k''\sqrt{m+n}$；\fld{thermal\_m/thermal\_n} 随结果返回；
  \item \textbf{剩余电压} \fld{v\_remaining\_pu} 与逐支路故障电流 \fld{SCDetailedBranchResult}（\srcpath{src/short\_circuit/short\_circuit.cpp:TransformerBranchCorrection}）。
\end{itemize}

\subsection{计算口径}
\fld{SCCalcType}（\fld{Max}/\fld{Min}，\srcpath{include/hacdcpf/analysis/short\_circuit.hpp:SCCalcType}）选择 $c_{\max}/c_{\min}$ 与相应机组数据；
\fld{SCTopology}（\fld{Auto}/\fld{Radial}/\fld{Meshed}，默认 \fld{Auto}）仅影响方法 B 的
网状 1.15 修正；低压 $\kappa$ 上限为 1.8，中高压为 2.0。\fld{Auto} 从实际参与电源到
故障点枚举独立路径，且任一路径 $R/X<0.3$ 时按标准不施加 1.15。开断电流 $I_b$ 的
公式 (77) 不由该拓扑枚举另行分支。
